\chapter*{Kurzfassung}
\addcontentsline{toc}{chapter}{Kurzfassung}
\markboth{Kurzfassung}{Kurzfassung}
\thispagestyle{plain}

Gestapelte mobile Bauzaunfelder sollen automatisiert gezählt werden -- über die
Detektion der Stirnflächen ihrer Rahmenrohre, kleiner und untereinander nahezu
identischer Kreisringe, die dicht gepackt und wechselseitig verdeckt auftreten.
Für dieses Zielobjekt existieren keine annotierten Datensätze, und die
Handannotation hunderter Kleinobjekte pro Bild ist teuer. Der hier verfolgte
Ansatz umgeht das Problem vollständig: Trainiert wird \emph{ausschließlich} auf
synthetischen Bildern, die aus neun freigestellten Objektausschnitten und sieben
Hintergrundaufnahmen erzeugt werden; validiert und getestet wird auf echten
Aufnahmen.

Damit verschiebt sich der methodische Schwerpunkt von der Netzarchitektur auf die
\textbf{Augmentierungs- und Synthese-Pipeline}: Sie ist hier nicht Beiwerk zu
einem vorhandenen Datensatz, sondern erzeugt die Trainingsverteilung vollständig.
Die Arbeit behandelt daher zunächst die theoretischen Grundlagen der
Datenaugmentierung -- Augmentierung als Kodierung von Invarianzen, als
Rauschregularisierung und als Mittel gegen \enquote{Shortcut Learning} -- und
entwickelt daraus ein Leitprinzip: Jede Transformation muss entweder eine an den
Realdaten \emph{gemessene} Eigenschaft nachbilden, eine \emph{physikalische
Randbedingung} respektieren oder eine konkrete \emph{Abkürzung} schließen, über
die das Netz sonst \enquote{hier wurde etwas eingeklebt} anstelle von \enquote{hier
ist ein Rohrende} lernen könnte. Die resultierende Pipeline umfasst 36
dokumentierte Transformationen; für jede werden Parameterbereich, gemessener
Ist-Anteil im erzeugten Datensatz (6000 Bilder, 173\,100 gerenderte Instanzen,
136\,384 gelabelte Objekte) und Begründung angegeben.

Auf zwei unberührten realen Aufnahmen zählt das Verfahren 86 statt 88 und 84
statt 86 Objekte -- ein \textbf{Zählfehler von 2,3\,\%}, die Zielgröße der
Aufgabenstellung. Die Bewertung liefert darüber hinaus fünf Befunde. Erstens ist
der Sim-to-Real-Abstand messbar und beträgt 9,6\PP{} \mAPfifty{}
(0,982 synthetisch gegen 0,886/0,888 real) -- klein genug, dass das Verfahren
funktioniert, und groß genug, um zu belegen, dass das Netz nicht bloß
Einfügeartefakte gelernt hat. Zweitens dominiert nicht eine
Augmentierungsentscheidung, sondern die Inferenzauflösung: \mAPfifty{} steigt von
0,435 bei 640\px{} auf 0,904 bei 1600\px. Drittens verbessert stärkere
Verdeckung in den Trainingsdaten das Ergebnis \emph{monoton} über zwei Splits und
drei Metriken -- obwohl der stärkste Arm 33\,\% weniger Labels enthält.
Viertens hebt ein zusätzlicher Detektionskopf auf Schrittweite~4 \mAPfifty{} auf
dem Testbild von 0,905 auf 0,941, und zwar -- wie aus der Zahl der
Kandidatengitterzellen je Objekt vorhergesagt -- gezielt dort, wo die Objekte
weniger als zwei Gitterzellen belegen; er verschlechtert dabei jedoch den
Zählfehler von 3,5 auf 8,1\,\%, weil sich das $F_1$-optimale
Konfidenzoptimum verschiebt. \mAPfifty{} und Zählfehler ordnen die Modelle also
gegenläufig -- für eine Zählanwendung ist mAP das falsche Auswahlkriterium.

Fünftens deckt die Fehleranalyse einen systematischen Größenversatz von 1,2 bis
2,5\px{} zwischen Generatorlabel und Handannotation auf. Er erklärt den auffällig
niedrigen \mAPseven{}-Wert auf dem Testbild -- eine Sensitivitätsrechnung hebt
ihn von 0,320 auf 0,714, ohne dass am Modell etwas geändert wird -- und lässt
\mAPfifty{}, $F_1$ und den Zählfehler unberührt. Bei programmatisch erzeugten
Labels wird die Labelkonvention damit selbst zu einer messbaren Größe.
